% ============================================================================
% 06_macros.tex  —  project macros
%
% THE CANONICAL-NUMBERS RULE (CLAUDE.md §5):
%   No numeric result may be hard-coded in any .tex file. Every number reaches
%   LaTeX through \num{key}, generated from data/CANONICAL_NUMBERS.yaml by
%   scripts/emit_numbers.py into report/preamble/numbers.tex.
%
% THE THREE PLACEHOLDERS (CLAUDE.md §4). Never silently fill any of them.
%   \PENDING{key}{what is needed}      a missing NUMBER
%   \NEEDSDATA{id}{which dataset}      a missing FIGURE or TABLE
%   \TODOPAL{question}                 something only Palaash can answer
% ============================================================================

\usepackage{xcolor}
\usepackage{tikz}
\usepackage{tcolorbox}
\tcbuselibrary{skins,breakable}

% ---------------------------------------------------------------------------
% \num{key} — canonical-number lookup
%
% numbers.tex defines one \DeclareNum{key}{rendered value} per VERIFIED entry.
% A key with no definition is by construction still PENDING, and renders as a
% loud red marker rather than silently producing nothing.
% ---------------------------------------------------------------------------
\makeatletter
\newcommand{\DeclareNum}[2]{\expandafter\gdef\csname yhsa@num@#1\endcsname{#2}}
\newcommand{\num}[1]{%
  \ifcsname yhsa@num@#1\endcsname
    \csname yhsa@num@#1\endcsname
  \else
    \yhsa@missingnum{#1}%
  \fi
}
\makeatother

% ---------------------------------------------------------------------------
% Semantic macros
% ---------------------------------------------------------------------------
\newcommand{\Pb}{\ce{Pb^2+}}
\newcommand{\Cu}{\ce{Cu^2+}}
\newcommand{\Zn}{\ce{Zn^2+}}
\newcommand{\PbII}{Pb(II)}
\newcommand{\CuII}{Cu(II)}
\newcommand{\ZnII}{Zn(II)}
\newcommand{\TAOSS}{TA--OSS}
\newcommand{\RAWOSS}{RAW--OSS}
\newcommand{\qmax}{q_{\mathrm{max}}}
\newcommand{\qe}{q_{\mathrm{e}}}
\newcommand{\Ce}{C_{\mathrm{e}}}
\newcommand{\Kd}{K_{\mathrm{d}}}
\newcommand{\forb}{f_{\mathrm{orb}}}
\newcommand{\BVten}{\mathrm{BV}_{10}}

% ---------------------------------------------------------------------------
% \finding{n}{text} — the numbered Findings (Bible §10, Device 3)
% Stated in bold at the point they are established, recalled by number in the
% Conclusion. Gives the referee a skeleton to score against.
% ---------------------------------------------------------------------------
\newcounter{yhsafinding}
\newcommand{\finding}[2]{%
  \par\medskip
  \begin{tcolorbox}[
    enhanced, breakable, sharp corners,
    colback=black!3, colframe=black!55, boxrule=0.8pt,
    left=8pt, right=8pt, top=6pt, bottom=6pt,
    before skip=6pt, after skip=10pt]
    \setstretch{1.0}\noindent\textbf{Finding #1.}\ #2
  \end{tcolorbox}
  \par\medskip
}

% ---------------------------------------------------------------------------
% \noveltybox — the boxed hypothesis / aim / contributions block (§1.4)
% Bible §10, Device 1. This is where a referee should be able to state the
% novelty in one sentence.
% ---------------------------------------------------------------------------
\newenvironment{noveltybox}[1][Hypothesis, aim and original contributions]{%
  \begin{tcolorbox}[
    enhanced, breakable, sharp corners,
    colback=black!2, colframe=black!70, boxrule=1pt,
    left=10pt, right=10pt, top=8pt, bottom=8pt,
    title={\normalfont\bfseries #1}, fonttitle=\bfseries,
    coltitle=white, colbacktitle=black!70]
  \setstretch{1.05}%
}{\end{tcolorbox}}

% ---------------------------------------------------------------------------
% \verdict{} — the mandatory closing sentence of every Results subsection
% Bible §8: "each subsection = one claim = one figure or table = one closing
% sentence that states what has now been established. Never leave a figure
% without a verdict."
% ---------------------------------------------------------------------------
\newcommand{\verdict}[1]{%
  \par\medskip
  \noindent\begin{minipage}{\linewidth}%
    \color{black!85}\rule{\linewidth}{0.4pt}\par\vspace{2pt}%
    \textit{#1}\par\vspace{2pt}%
    \rule{\linewidth}{0.4pt}%
  \end{minipage}\par\medskip
}

% ---------------------------------------------------------------------------
% \origin{owner}{text} — origin tagging (Bible §10, Device 2)
% Names the owner of every existing method in the sentence that uses it, so the
% distinction between background and original contribution is mechanical rather
% than implied. Criterion C7 asks for this literally.
% ---------------------------------------------------------------------------
\newcommand{\origin}[2]{#2\,\textsuperscript{\textit{[#1]}}}
\newcommand{\ourwork}[1]{#1}

% ---------------------------------------------------------------------------
% Placeholders. Rendering is set by 07_draft_guards.tex: visible red boxes in
% draft mode, hard build failure in final mode.
% ---------------------------------------------------------------------------
% \makeatletter is required: the internal names contain @. Without it,
% \yhsa@pending parses as \yhsa followed by literal '@pending'.
\makeatletter
\newcommand{\PENDING}[2]{\yhsa@pending{#1}{#2}}
\newcommand{\NEEDSDATA}[2]{\yhsa@needsdata{#1}{#2}}
\newcommand{\TODOPAL}[1]{\yhsa@todopal{#1}}
\makeatother
